В качестве базового средства распознавания текста в работе использовалась библиотека EasyOCR. Данное программное средство представляет собой готовую OCR-систему с открытым исходным кодом, ориентированную на распознавание текста на изображениях и сканах документов. Библиотека поддерживает большое число языков, включая русский, и может быть легко интегрирована в программный конвейер на языке Python. EasyOCR удобна тем, что объединяет два ключевых этапа OCR-обработки: обнаружение текстовых областей на изображении и последующее распознавание символов в найденных фрагментах. Это позволяет использовать ее как базовый модуль распознавания без необходимости самостоятельно реализовывать этап сегментации текста.

Вместе с тем качество работы EasyOCR существенно зависит от состояния входного изображения. При наличии плохого качества сканированной книге на входе число ошибок распознавания возрастает. По этой причине в работе библиотека рассматривалась не только как самостоятельное средство OCR, но и как основа для дальнейшего улучшения результатов с помощью адаптивной предобработки изображения и последующей нормализации распознанного текста.